\documentclass[10pt]{article}
\usepackage[margin=0.85in]{geometry}
\usepackage{amsmath,amssymb}
\usepackage[hidelinks]{hyperref}
\usepackage{microtype}

\title{Literature Status: Unique Preduals for\\
Finite Maximal Subdiagonal Algebras}
\author{Run \texttt{fa\_banach\_001} --- lane 13}
\date{9 August 2026}

\begin{document}
\maketitle

\noindent\textbf{Status.} \texttt{literature\_already\_answered}.

\section*{Source question}

In Closing Remark 2 on PDF page 21 of \cite{BL2007}, Blecher and
Labuschagne write:
\begin{quote}
A nice question communicated to us by Gilles Godefroy at this conference is
whether subdiagonal algebras have unique preduals.
\end{quote}
Their standing setting is a finite von Neumann algebra $M$ with faithful normal
tracial state $\tau$ and a finite maximal subdiagonal algebra
$A=H^\infty(M,\tau)$.

\section*{Exact later answer}

Ueda's introduction explicitly says that his unique-predual result is ``an
affirmative answer to a question posed by Godefroy stated in'' the
Blecher--Labuschagne survey \cite[p.~2]{Ueda2008}.  His Theorem 3.1, on PDF
page 4, states
\[
       M_*/A_\perp \text{ is the unique predual of }
       A=H^\infty(M,\tau).
\]
This is precisely the canonical predual and exactly the class of algebras in
the source question.  Thus the answer is yes.

\section*{Provenance and scope}

The answering author knew he was resolving this exact survey question, so this
is an explicit literature answer rather than an agent-derived implication.
The packet does not classify preduals outside the finite tracial maximal
subdiagonal setting and does not address the other questions in the survey's
closing remarks.

\begin{thebibliography}{2}

\bibitem{BL2007}
D. P. Blecher and L. E. Labuschagne,
\emph{Von Neumann algebraic $H^p$ theory},
arXiv:math/0611879 (2006/2007).

\bibitem{Ueda2008}
Y. Ueda,
\emph{On peak phenomena for non-commutative $H^\infty$},
arXiv:0802.3449; Math. Ann. \textbf{343} (2009), 421--429.

\end{thebibliography}

\end{document}
